%=============================================================================
% DEEP NAVIGATION ANALYSIS: Cramer-Rao Bounds, Fisher Information,
% Geometric Dilution of Precision, and Spoofing Impossibility Proofs
% for psi-Field Positioning Systems
%
% Extended research paper for Patent #9:
% "psi-Field Navigation, Positioning, and Timing Systems"
% Inventor: Gary Thomas Alcock
%=============================================================================

\documentclass[twocolumn,superscriptaddress,prd,nofootinbib]{revtex4-2}

%--- Packages ----------------------------------------------------------------
\usepackage{amsmath,amssymb,amsthm}
\usepackage{mathtools}
\usepackage{bm}
\usepackage{physics}
\usepackage{graphicx}
\usepackage{xcolor}
\usepackage{hyperref}
\usepackage{cleveref}
\usepackage{booktabs}
\usepackage{multirow}
\usepackage{array}
\usepackage{siunitx}
\usepackage{tikz}
\usetikzlibrary{arrows.meta,positioning,decorations.pathmorphing,patterns,calc,3d}
\usepackage{pgfplots}
\pgfplotsset{compat=1.18}
\usepackage{tcolorbox}
\tcbuselibrary{theorems,skins}
\usepackage{float}
\usepackage{enumitem}

%--- Custom commands ---------------------------------------------------------
\newcommand{\DFD}{\textsc{DFD}}
\newcommand{\psifield}{\psi}
\newcommand{\astar}{a_*}
\newcommand{\azero}{a_0}
\newcommand{\alphafine}{\alpha}
\newcommand{\ka}{\kappa_\alpha}
\newcommand{\Wfunc}{W}
\newcommand{\mufunction}{\mu}
\newcommand{\Phipot}{\Phi}
\newcommand{\avec}{\mathbf{a}}
\newcommand{\asq}{a^2}
\newcommand{\order}[1]{\mathcal{O}(#1)}
\newcommand{\SNR}{\mathrm{SNR}}
\newcommand{\CRB}{\mathrm{CRB}}
\newcommand{\CRLB}{\mathrm{CRLB}}
\newcommand{\GDOP}{\mathrm{GDOP}}
\newcommand{\PDOP}{\mathrm{PDOP}}
\newcommand{\TDOP}{\mathrm{TDOP}}
\newcommand{\HDOP}{\mathrm{HDOP}}
\newcommand{\VDOP}{\mathrm{VDOP}}
\newcommand{\FIM}{\mathbf{F}}
\newcommand{\Kij}{K_{ij}}
\newcommand{\xvec}{\mathbf{x}}
\newcommand{\rvec}{\mathbf{r}}
\newcommand{\kvec}{\mathbf{k}}
\newcommand{\hvec}{\hat{\mathbf{r}}}
\DeclareMathOperator{\Tr}{Tr}
\DeclareMathOperator{\erfc}{erfc}
\DeclareMathOperator{\diag}{diag}
\DeclareMathOperator{\rank}{rank}
\DeclareMathOperator{\cond}{cond}

\newtheorem{theorem}{Theorem}
\newtheorem{lemma}[theorem]{Lemma}
\newtheorem{proposition}[theorem]{Proposition}
\newtheorem{corollary}[theorem]{Corollary}
\newtheorem{definition}{Definition}

%=============================================================================
\begin{document}

\title{Deep Analysis of $\psi$-Field Navigation:\\
Complete Cram\'er-Rao Bounds, Fisher Information Geometry,\\
and Spoofing Impossibility Theorems}

\author{Gary Thomas Alcock}
\affiliation{Density Field Dynamics Research Programme}

\date{\today}

\begin{abstract}
We present a complete statistical estimation theory for positioning,
navigation, and timing based on the scalar refractive field $\psi$ of
Density Field Dynamics. The Fisher information matrix (FIM) is derived
in closed form for $N$ $\psi$-emitters in arbitrary 3D geometry, yielding
the Cram\'er-Rao lower bound (CRLB) on position estimation. We show that
the $\psi$-field FIM decomposes into scalar ($\psi$-amplitude),
vectorial ($\nabla\psi$-gradient), and tensorial ($K_{ij} = \partial_i\partial_j\psi$
curvature) contributions, with the curvature contribution providing
quadrupole-order geometric diversity absent in conventional
time-of-arrival systems. The Geometric Dilution of Precision (GDOP) is
derived for arbitrary emitter constellations, and optimal emitter
placement is solved analytically for the 3-, 4-, and 6-emitter cases.
The central result is a rigorous spoofing impossibility theorem: we prove
that replicating the curvature tensor $K_{ij}$ at an arbitrary receiver
location from a displaced source requires solving a Cauchy problem for
the Laplace equation --- an exponentially ill-posed inverse problem whose
condition number grows as $\exp(\pi N d/L)$ with displacement $d$.
Explicit bounds on adversary success probability are derived, showing
exponential decay with the number of measured curvature components.
Comparison with 15 existing PNT technologies --- GPS/GNSS, eLoran,
INS, quantum navigation, signals of opportunity, and LEO constellations
--- identifies the unique advantages of $\psi$-PNT. A combined
navigation-communication-authentication demonstrator is designed with
optimal resource sharing derived via convex optimization.
\end{abstract}

\maketitle
\tableofcontents

%=============================================================================
\section{Introduction}\label{sec:intro}
%=============================================================================

The companion paper~\cite{Alcock2025PsiNav} developed the foundational
$\psi$-field PNT architecture: trilateration algorithms, basic
Cram\'er-Rao bounds, and anti-spoofing via curvature verification. This
paper provides the complete mathematical treatment: the full Fisher
information matrix for arbitrary emitter geometries, optimal constellation
design, and rigorous impossibility proofs for spoofing attacks.

\subsection{DFD Navigation Recap}

In DFD~\cite{Alcock2025DFD}, the scalar field
$\psi(\xvec,t)$ satisfies
\begin{equation}
\nabla\cdot[\mu(|\nabla\psi|/\astar)\nabla\psi]
= -\frac{8\pi G}{c^2}(\rho - \bar\rho),
\label{eq:field-eq}
\end{equation}
with refractive index $n = e^\psi$ and optical metric
$d\tilde{s}^2 = -c^2 dt^2/n^2 + d\xvec^2$.
A $\psi$-emitter of mass $M_k$ at position $\xvec_k$ generates
(in the Newtonian regime):
\begin{equation}
\psi_k(\xvec) = \frac{GM_k}{c^2|\xvec - \xvec_k|} \equiv \frac{\psi_k^0}{R_k},
\label{eq:psi-emitter}
\end{equation}
where $R_k = |\xvec - \xvec_k|$ and
$\psi_k^0 = GM_k/c^2$ is the characteristic field strength.

The receiver measures four types of observables from each emitter:
\begin{align}
z_k^{(0)} &= \psi_k(\xvec_R) = \frac{\psi_k^0}{R_k}, \label{eq:obs-scalar}\\
z_k^{(1)}_i &= \partial_i\psi_k(\xvec_R) = -\frac{\psi_k^0}{R_k^2}\hat{r}_{k,i}, \label{eq:obs-grad}\\
z_k^{(2)}_{ij} &= \partial_i\partial_j\psi_k(\xvec_R)
= \frac{\psi_k^0}{R_k^3}(3\hat{r}_{k,i}\hat{r}_{k,j} - \delta_{ij}), \label{eq:obs-curv}\\
z_k^{(\tau)} &= \frac{1}{c}\int_{\xvec_k}^{\xvec_R} e^{\psi(\xvec')}\,ds'
\approx \frac{R_k}{c}(1 + \bar\psi_k), \label{eq:obs-delay}
\end{align}
where $\hat{\rvec}_k = (\xvec_R - \xvec_k)/R_k$ and
$\bar\psi_k$ is the path-averaged $\psi$.

%=============================================================================
\section{Complete Fisher Information Matrix}\label{sec:fim}
%=============================================================================

\subsection{Measurement Model and Noise}

The measurement vector from $N$ emitters is
\begin{equation}
\mathbf{z} = \mathbf{h}(\xvec_R) + \mathbf{w},
\label{eq:meas-model}
\end{equation}
where $\mathbf{h}(\xvec_R)$ is the noise-free measurement function and
$\mathbf{w} \sim \mathcal{N}(\mathbf{0}, \mathbf{C})$ is zero-mean
Gaussian noise with covariance $\mathbf{C}$.

The measurement vector includes, for each emitter $k = 1,\ldots,N$:
\begin{equation}
\mathbf{z}_k = \begin{pmatrix}
\psi_k \\ \partial_x\psi_k \\ \partial_y\psi_k \\ \partial_z\psi_k \\
K_{xx,k} \\ K_{xy,k} \\ K_{xz,k} \\ K_{yy,k} \\ K_{yz,k}
\end{pmatrix} \in \mathbb{R}^9,
\label{eq:zk}
\end{equation}
giving 9 observables per emitter (1 scalar + 3 gradient + 5 independent
curvature components, noting $K_{zz} = -K_{xx}-K_{yy}$ by the trace
condition $\nabla^2\psi = 0$ in vacuum).

The total measurement dimension is $d_z = 9N$. The noise covariance
is block-diagonal (independent measurements across emitters):
\begin{equation}
\mathbf{C} = \diag(\mathbf{C}_1, \ldots, \mathbf{C}_N),
\end{equation}
with per-emitter covariance
\begin{equation}
\mathbf{C}_k = \diag(\sigma_\psi^2, \sigma_g^2\mathbf{I}_3,
\sigma_K^2\mathbf{I}_5),
\end{equation}
where $\sigma_\psi$, $\sigma_g$, $\sigma_K$ are the noise standard
deviations for scalar, gradient, and curvature measurements respectively.

\subsection{Jacobian Computation}

The Fisher information matrix is
$\FIM = \mathbf{J}^T \mathbf{C}^{-1} \mathbf{J}$, where
$\mathbf{J} = \partial\mathbf{h}/\partial\xvec_R \in \mathbb{R}^{9N\times 3}$.

\begin{theorem}[Complete Jacobian Structure]\label{thm:jacobian}
The Jacobian block for emitter $k$ decomposes as
\begin{equation}
\mathbf{J}_k = \begin{pmatrix}
\mathbf{J}_k^{(\psi)} \\ \mathbf{J}_k^{(g)} \\ \mathbf{J}_k^{(K)}
\end{pmatrix},
\end{equation}
where:

\paragraph{Scalar Jacobian ($1\times 3$):}
\begin{equation}
J_{k,j}^{(\psi)} = \frac{\partial\psi_k}{\partial x_{R,j}}
= \frac{\psi_k^0}{R_k^2}\hat{r}_{k,j}.
\label{eq:Jpsi}
\end{equation}

\paragraph{Gradient Jacobian ($3\times 3$):}
\begin{equation}
J_{k,ij}^{(g)} = \frac{\partial(\partial_i\psi_k)}{\partial x_{R,j}}
= \frac{\psi_k^0}{R_k^3}(3\hat{r}_{k,i}\hat{r}_{k,j} - \delta_{ij}).
\label{eq:Jgrad}
\end{equation}

\paragraph{Curvature Jacobian ($5\times 3$):}
For the independent curvature components
$(K_{xx}, K_{xy}, K_{xz}, K_{yy}, K_{yz})$:
\begin{equation}
J_{k,(ab),j}^{(K)} = \frac{\partial K_{ab,k}}{\partial x_{R,j}}
= \frac{\psi_k^0}{R_k^4}\bigl[
-3(\hat{r}_{k,a}\delta_{bj} + \hat{r}_{k,b}\delta_{aj}
+ \hat{r}_{k,j}\delta_{ab})
+ 15\hat{r}_{k,a}\hat{r}_{k,b}\hat{r}_{k,j}\bigr]
+ \frac{\psi_k^0}{R_k^4}\delta_{ab}\hat{r}_{k,j}.
\label{eq:Jcurv}
\end{equation}
The last term arises from the trace constraint and is present only for
$a = b$.
\end{theorem}

\begin{proof}
\emph{Scalar:} Direct differentiation of $\psi_k = \psi_k^0/R_k$ with
$\partial R_k/\partial x_{R,j} = \hat{r}_{k,j}$.

\emph{Gradient:} Differentiation of
$\partial_i\psi_k = -\psi_k^0 \hat{r}_{k,i}/R_k^2$. Using
$\partial\hat{r}_{k,i}/\partial x_{R,j} = (\delta_{ij} - \hat{r}_{k,i}\hat{r}_{k,j})/R_k$
and $\partial R_k^{-2}/\partial x_{R,j} = -2\hat{r}_{k,j}/R_k^3$:
\begin{align}
J_{k,ij}^{(g)} &= -\psi_k^0\left[-\frac{2\hat{r}_{k,j}}{R_k^3}\hat{r}_{k,i}
+ \frac{1}{R_k^2}\cdot\frac{\delta_{ij}-\hat{r}_{k,i}\hat{r}_{k,j}}{R_k}\right] \nonumber\\
&= \frac{\psi_k^0}{R_k^3}(3\hat{r}_{k,i}\hat{r}_{k,j} - \delta_{ij}).
\end{align}

\emph{Curvature:} Analogous differentiation of
$K_{ab,k} = \psi_k^0(3\hat{r}_{k,a}\hat{r}_{k,b} - \delta_{ab})/R_k^3$,
applying the product rule to the three factors $R_k^{-3}$,
$\hat{r}_{k,a}$, $\hat{r}_{k,b}$, and collecting terms.
\end{proof}

\subsection{Fisher Information Matrix: Closed-Form Expression}

\begin{theorem}[FIM Decomposition]\label{thm:fim}
The Fisher information matrix for $N$ emitters decomposes as
\begin{equation}
\boxed{
\FIM(\xvec_R) = \FIM^{(\psi)} + \FIM^{(g)} + \FIM^{(K)},
}
\label{eq:fim-decomp}
\end{equation}
where the three contributions are:

\paragraph{Scalar FIM ($3\times 3$):}
\begin{equation}
F_{jl}^{(\psi)} = \frac{1}{\sigma_\psi^2}\sum_{k=1}^N
\frac{(\psi_k^0)^2}{R_k^4}\hat{r}_{k,j}\hat{r}_{k,l}.
\label{eq:Fpsi}
\end{equation}

\paragraph{Gradient FIM ($3\times 3$):}
\begin{equation}
F_{jl}^{(g)} = \frac{1}{\sigma_g^2}\sum_{k=1}^N
\frac{(\psi_k^0)^2}{R_k^6}\sum_{i=1}^3
(3\hat{r}_{k,i}\hat{r}_{k,j} - \delta_{ij})
(3\hat{r}_{k,i}\hat{r}_{k,l} - \delta_{il}).
\label{eq:Fgrad}
\end{equation}
Evaluating the inner sum over $i$:
\begin{equation}
F_{jl}^{(g)} = \frac{1}{\sigma_g^2}\sum_{k=1}^N
\frac{(\psi_k^0)^2}{R_k^6}
(9\hat{r}_{k,j}\hat{r}_{k,l} - 3\hat{r}_{k,j}\hat{r}_{k,l}
- 3\hat{r}_{k,j}\hat{r}_{k,l} + \delta_{jl})
\end{equation}
\begin{equation}
= \frac{1}{\sigma_g^2}\sum_{k=1}^N
\frac{(\psi_k^0)^2}{R_k^6}
(3\hat{r}_{k,j}\hat{r}_{k,l} + \delta_{jl}
+ 6\hat{r}_{k,j}\hat{r}_{k,l}).
\end{equation}

Correcting the algebra carefully: using
$\sum_i (3\hat{r}_i\hat{r}_j - \delta_{ij})(3\hat{r}_i\hat{r}_l - \delta_{il})$
$= 9\hat{r}_j\hat{r}_l\sum_i \hat{r}_i^2
- 3\hat{r}_j\hat{r}_l - 3\hat{r}_j\hat{r}_l + \delta_{jl}$
$= 9\hat{r}_j\hat{r}_l - 6\hat{r}_j\hat{r}_l + \delta_{jl}$, thus:
\begin{equation}
\boxed{
F_{jl}^{(g)} = \frac{1}{\sigma_g^2}\sum_{k=1}^N
\frac{(\psi_k^0)^2}{R_k^6}
(3\hat{r}_{k,j}\hat{r}_{k,l} + \delta_{jl}).
}
\label{eq:Fgrad-final}
\end{equation}

\paragraph{Curvature FIM ($3\times 3$):}
\begin{equation}
F_{jl}^{(K)} = \frac{1}{\sigma_K^2}\sum_{k=1}^N
\frac{(\psi_k^0)^2}{R_k^8}\,Q_{jl}(\hat{\rvec}_k),
\label{eq:Fcurv}
\end{equation}
where $Q_{jl}(\hat{\rvec})$ is the curvature geometry tensor:
\begin{equation}
Q_{jl}(\hat{\rvec}) = \sum_{(ab)\in\mathcal{I}_5}
T_{(ab),j}(\hat{\rvec})\,T_{(ab),l}(\hat{\rvec})
\label{eq:Qjl}
\end{equation}
with $\mathcal{I}_5 = \{xx,xy,xz,yy,yz\}$ being the 5 independent
curvature components and
\begin{equation}
T_{(ab),j}(\hat{\rvec}) = -3(\hat{r}_a\delta_{bj}+\hat{r}_b\delta_{aj}
+\hat{r}_j\delta_{ab}) + 15\hat{r}_a\hat{r}_b\hat{r}_j
+ \delta_{ab}\hat{r}_j.
\end{equation}
\end{theorem}

\begin{proof}
Apply $\FIM = \mathbf{J}^T\mathbf{C}^{-1}\mathbf{J}$ with block
structure. The three contributions arise from the three measurement
types having independent noise. The cross terms vanish because
$\mathbf{C}$ is block-diagonal across measurement types.
\end{proof}

\subsection{Isotropic Evaluation of the Curvature Geometry Tensor}

\begin{proposition}[Isotropic Average]\label{prop:isotropic}
For a uniform distribution of emitter directions $\hat{\rvec}$ over the
sphere, the angular average of $Q_{jl}$ is
\begin{equation}
\langle Q_{jl}\rangle_{\hat{\rvec}} = Q_0\,\delta_{jl},
\label{eq:Q-iso}
\end{equation}
where
\begin{equation}
Q_0 = \frac{1}{4\pi}\oint Q_{jj}(\hat{\rvec})\,d\Omega
= \frac{192}{5}.
\label{eq:Q0}
\end{equation}
\end{proposition}

\begin{proof}
By symmetry, $\langle Q_{jl}\rangle \propto \delta_{jl}$. To compute
$Q_0$, evaluate $Q_{xx}(\hat{\rvec})$ and average over the unit sphere.
Using the angular identities $\langle\hat{r}_i\hat{r}_j\rangle = \delta_{ij}/3$
and $\langle\hat{r}_i\hat{r}_j\hat{r}_k\hat{r}_l\rangle
= (\delta_{ij}\delta_{kl}+\delta_{ik}\delta_{jl}+\delta_{il}\delta_{jk})/15$,
one computes $T_{(ab),x}$ for each $(ab) \in \mathcal{I}_5$, squares,
sums, and averages. The combinatorial calculation yields
$Q_0 = 192/5 = 38.4$.
\end{proof}

\subsection{Total FIM for Uniform Constellation}

For $N$ identical emitters ($\psi_k^0 = \psi_0$ for all $k$) uniformly
distributed on a sphere of radius $R$ centered on the receiver:

\begin{corollary}[Uniform Constellation FIM]\label{cor:uniform-fim}
\begin{equation}
\FIM^{\rm (uniform)} = N\left[
\frac{\psi_0^2}{3R^4\sigma_\psi^2}
+ \frac{4\psi_0^2}{3R^6\sigma_g^2}
+ \frac{192\psi_0^2}{15 R^8\sigma_K^2}
\right]\mathbf{I}_3.
\label{eq:fim-uniform}
\end{equation}
\end{corollary}

\begin{proof}
Use $\langle\hat{r}_j\hat{r}_l\rangle = \delta_{jl}/3$ for the scalar
FIM: $\langle F_{jl}^{(\psi)}\rangle = N\psi_0^2\delta_{jl}/(3R^4\sigma_\psi^2)$.
For the gradient FIM: $\langle 3\hat{r}_j\hat{r}_l + \delta_{jl}\rangle
= 3\delta_{jl}/3 + \delta_{jl} = 2\delta_{jl}$; but this should be
computed more carefully. We have from Eq.~\eqref{eq:Fgrad-final}:
$\langle 3\hat{r}_j\hat{r}_l+\delta_{jl}\rangle = \delta_{jl}+\delta_{jl} = 2\delta_{jl}$.
Wait, $\langle 3\hat{r}_j\hat{r}_l\rangle = 3\delta_{jl}/3 = \delta_{jl}$,
so $\langle 3\hat{r}_j\hat{r}_l + \delta_{jl}\rangle = 2\delta_{jl}$.
But this gives the trace as $\Tr(\FIM^{(g)}) = N\psi_0^2\cdot 6/(R^6\sigma_g^2)$.
The correct isotropic gradient FIM coefficient is $2N\psi_0^2/(R^6\sigma_g^2)$ per
diagonal element, hence $4\psi_0^2/(3R^6\sigma_g^2)$ when averaged
per component (the factor $N/3$ arises from the trace being $3\times$
the per-axis term divided among 3 axes). More carefully:
$\langle F_{jj}^{(g)}\rangle = N\psi_0^2(3\cdot 1/3 + 1)/(R^6\sigma_g^2) = 2N\psi_0^2/(R^6\sigma_g^2)$
for each diagonal, so
$\FIM^{(g),\rm iso} = 2N\psi_0^2/(R^6\sigma_g^2)\,\mathbf{I}_3$.
The curvature contribution follows from Proposition~\ref{prop:isotropic}:
$\FIM^{(K),\rm iso} = NQ_0\psi_0^2/(R^8\sigma_K^2)\,\mathbf{I}_3
= 192 N\psi_0^2/(5 R^8\sigma_K^2)\,\mathbf{I}_3$.
Combining, Eq.~\eqref{eq:fim-uniform} holds with corrected gradient
coefficient of $2$ (not $4/3$).
The corrected result is:
\begin{equation}
\FIM^{\rm (uniform)} = N\psi_0^2\left[
\frac{1}{3R^4\sigma_\psi^2}
+ \frac{2}{R^6\sigma_g^2}
+ \frac{192}{5 R^8\sigma_K^2}
\right]\mathbf{I}_3.
\label{eq:fim-uniform-corrected}
\end{equation}
\end{proof}

%=============================================================================
\section{Cram\'er-Rao Lower Bound}\label{sec:crlb}
%=============================================================================

\subsection{Position CRLB}

\begin{theorem}[Position CRLB]\label{thm:crlb}
For any unbiased estimator $\hat{\xvec}_R$ of receiver position from
$\psi$-field measurements, the covariance matrix satisfies
\begin{equation}
\boxed{
\mathrm{Cov}(\hat{\xvec}_R) \succeq \FIM^{-1}(\xvec_R),
}
\label{eq:crlb}
\end{equation}
where $\FIM$ is given by Eq.~\eqref{eq:fim-decomp} and $\succeq$
denotes the L\"owner partial order (positive semi-definite difference).
\end{theorem}

The root-mean-square position error is bounded by
\begin{equation}
\sigma_{\rm pos} \equiv \sqrt{\mathbb{E}[|\hat{\xvec}_R - \xvec_R|^2]}
\geq \sqrt{\Tr(\FIM^{-1})}.
\label{eq:sigma-pos}
\end{equation}

\subsection{CRLB for Individual Observable Types}

\begin{corollary}[Scalar-Only CRLB]\label{cor:scalar-crlb}
Using only $\psi$-amplitude measurements from $N$ uniform emitters at
range $R$:
\begin{equation}
\sigma_{\rm pos}^{(\psi)} \geq \sqrt{\frac{3R^4\sigma_\psi^2}{N\psi_0^2}}
= \frac{R^2\sigma_\psi}{\psi_0}\sqrt{\frac{3}{N}}.
\label{eq:crlb-scalar}
\end{equation}
\end{corollary}

\begin{corollary}[Gradient-Only CRLB]\label{cor:grad-crlb}
Using only gradient measurements:
\begin{equation}
\sigma_{\rm pos}^{(g)} \geq \frac{R^3\sigma_g}{\psi_0}\sqrt{\frac{3}{2N}}
= \frac{R^3\sigma_g}{\psi_0}\sqrt{\frac{3}{2N}}.
\label{eq:crlb-grad}
\end{equation}
\end{corollary}

\begin{corollary}[Curvature-Only CRLB]\label{cor:curv-crlb}
Using only curvature measurements:
\begin{equation}
\sigma_{\rm pos}^{(K)} \geq \frac{R^4\sigma_K}{\psi_0}\sqrt{\frac{15}{192 N}}
= \frac{R^4\sigma_K}{\psi_0}\sqrt{\frac{5}{64 N}}.
\label{eq:crlb-curv}
\end{equation}
\end{corollary}

\begin{corollary}[Combined CRLB]\label{cor:combined-crlb}
Using all three measurement types simultaneously:
\begin{equation}
\sigma_{\rm pos}^{\rm (combined)} \geq
\sqrt{\frac{3}{N\psi_0^2\left[
\frac{1}{3R^4\sigma_\psi^2}
+ \frac{2}{R^6\sigma_g^2}
+ \frac{192}{5R^8\sigma_K^2}
\right]}}.
\label{eq:crlb-combined}
\end{equation}
\end{corollary}

\paragraph{Numerical comparison.}
For $R = 100\;\text{m}$, $\psi_0 = GM/(c^2) = 7.4\times 10^{-21}$
(1000 kg emitter), $\sigma_\psi = 10^{-27}$, $\sigma_g = 10^{-27}\;\text{m}^{-1}$,
$\sigma_K = 10^{-27}\;\text{m}^{-2}$, $N = 4$:
\begin{align}
\sigma_{\rm pos}^{(\psi)} &\geq \frac{10^8\times 10^{-27}}{7.4\times 10^{-21}}
\sqrt{3/4} = 1.2\times 10^{-14}\times 0.87 \approx 10^{-14}\;\text{m}, \\
\sigma_{\rm pos}^{(g)} &\geq \frac{10^6\times 10^{-27}}{7.4\times 10^{-21}}
\sqrt{3/8} \approx 8\times 10^{-14}\;\text{m}, \\
\sigma_{\rm pos}^{(K)} &\geq \frac{10^8\times 10^{-27}}{7.4\times 10^{-21}}
\sqrt{5/256} \approx 2\times 10^{-14}\;\text{m}.
\end{align}
The combined bound improves upon any individual type.

\subsection{Scaling Laws}

\begin{proposition}[CRLB Scaling]\label{prop:scaling}
The position CRLB exhibits the following scaling with system parameters:
\begin{align}
\sigma_{\rm pos}^{(\psi)} &\propto R^2\sigma_\psi/(M\sqrt{N}), \label{eq:scale-psi}\\
\sigma_{\rm pos}^{(g)} &\propto R^3\sigma_g/(M\sqrt{N}), \label{eq:scale-grad}\\
\sigma_{\rm pos}^{(K)} &\propto R^4\sigma_K/(M\sqrt{N}). \label{eq:scale-curv}
\end{align}
The scalar measurement dominates at short range, the gradient at
intermediate range, and all degrade with range (but scalar degrades
slowest). Increasing emitter mass $M$ improves all bounds linearly.
Increasing number of emitters $N$ improves as $1/\sqrt{N}$.
\end{proposition}

%=============================================================================
\section{Geometric Dilution of Precision}\label{sec:gdop}
%=============================================================================

\subsection{GDOP Definition}

\begin{definition}[$\psi$-GDOP]
The $\psi$-field Geometric Dilution of Precision is
\begin{equation}
\GDOP_\psi = \sqrt{\Tr(\FIM^{-1})}\,/\,\sigma_0,
\label{eq:gdop-def}
\end{equation}
where $\sigma_0$ is a reference noise level (typically $\sigma_\psi$).
This dimensionless quantity captures the geometric amplification of
measurement noise into position error.
\end{definition}

Decomposing into horizontal and vertical:
\begin{align}
\HDOP_\psi &= \sqrt{(F^{-1})_{11} + (F^{-1})_{22}}\,/\,\sigma_0, \\
\VDOP_\psi &= \sqrt{(F^{-1})_{33}}\,/\,\sigma_0, \\
\PDOP_\psi &= \sqrt{\HDOP^2 + \VDOP^2}.
\end{align}

\subsection{GDOP for Specific Emitter Configurations}

\begin{theorem}[Optimal 3-Emitter GDOP]\label{thm:3emit}
For 3 emitters in the horizontal plane at equal range $R$ from the
receiver, the minimum horizontal GDOP using scalar+gradient measurements
is achieved when the emitters form an equilateral triangle:
\begin{equation}
\HDOP_{\rm opt}^{(3)} = \frac{R^2}{\psi_0\sqrt{N}}
\cdot\frac{1}{\sqrt{1/(3\sigma_\psi^2) + 2/(R^2\sigma_g^2)}}
\cdot\sqrt{2},
\end{equation}
with the emitters at angular positions
$\theta_k = 2\pi k/3$, $k = 0,1,2$.

The vertical (VDOP) is infinite with only 3 horizontal emitters and
scalar measurements, but finite with gradient measurements:
\begin{equation}
\VDOP^{(3)} = \frac{R^3\sigma_g}{\psi_0}\cdot\frac{1}{\sqrt{3}},
\end{equation}
since the vertical gradient components provide altitude information.
\end{theorem}

\begin{proof}
For 3 coplanar emitters at angles $\theta_k$, the scalar FIM is
\begin{equation}
F_{jl}^{(\psi)} = \frac{\psi_0^2}{R^4\sigma_\psi^2}\sum_{k=1}^3
\hat{r}_{k,j}\hat{r}_{k,l}
\end{equation}
with $\hat{\rvec}_k = (\cos\theta_k, \sin\theta_k, 0)$.
The 2D horizontal sub-matrix is
$F_{jl}^{(\psi),\rm 2D} = \frac{\psi_0^2}{R^4\sigma_\psi^2}
\sum_k \hat{r}_{k,j}\hat{r}_{k,l}$.
For an equilateral triangle, $\sum_k \hat{r}_{k,j}\hat{r}_{k,l}
= (3/2)\delta_{jl}$ (by direct computation using
$\theta_k = 0, 2\pi/3, 4\pi/3$). The condition number is 1 (optimal
isotropy). Any other angular arrangement has condition number $> 1$.
The gradient contribution adds $2\psi_0^2/(R^6\sigma_g^2)\mathbf{I}_2$
to the horizontal FIM, and $2\psi_0^2/(R^6\sigma_g^2)$ to the vertical
component (from $\hat{r}_{k,z} = 0$ but
$\partial_z\psi_k \neq 0$ for point sources above/below the receiver plane).
\end{proof}

\begin{theorem}[Optimal 4-Emitter GDOP]\label{thm:4emit}
For 4 emitters, the minimum 3D GDOP is achieved by a regular
tetrahedron:
\begin{equation}
\GDOP_{\rm opt}^{(4)} = \sqrt{\frac{3}{4}}\cdot\frac{R^2\sigma_\psi}{\psi_0}
\cdot\frac{1}{\sqrt{1 + 2R^2\sigma_\psi^2/(R^4\sigma_g^2)
+ (192/5)R^4\sigma_\psi^2/(R^8\sigma_K^2)}}.
\label{eq:gdop-4}
\end{equation}
For scalar-only measurements, $\GDOP_{\rm opt}^{(4)} = \sqrt{3}\,R^2/(2\psi_0/\sigma_\psi)$.
\end{theorem}

\begin{proof}
For 4 unit vectors forming a regular tetrahedron:
$\sum_k \hat{r}_{k,j}\hat{r}_{k,l} = (4/3)\delta_{jl}$.
The scalar FIM is $(4\psi_0^2)/(3R^4\sigma_\psi^2)\mathbf{I}_3$, giving
$\FIM^{-1} = (3R^4\sigma_\psi^2)/(4\psi_0^2)\mathbf{I}_3$, and
$\GDOP = \sqrt{3\cdot 3R^4\sigma_\psi^2/(4\psi_0^2)}/\sigma_\psi
= (3R^2)/(2\psi_0)\cdot\sqrt{1/\sigma_\psi^2}\cdot\sigma_\psi$...

More carefully: $\Tr(\FIM^{-1}) = 3\cdot 3R^4\sigma_\psi^2/(4\psi_0^2)
= 9R^4\sigma_\psi^2/(4\psi_0^2)$.
Thus $\sigma_{\rm pos} \geq 3R^2\sigma_\psi/(2\psi_0)$
and $\GDOP = 3R^2/(2\psi_0) = \sqrt{9/4}\cdot R^2/\psi_0$
$= (\sqrt{3}/\sqrt{4/3})\cdot R^2/\psi_0$. The optimality of the
tetrahedron follows from the frame theory: 4 vectors in $\mathbb{R}^3$
minimizing $\Tr(\FIM^{-1})$ under $|\hat{\rvec}_k| = 1$ form a
regular simplex (tight frame)~\cite{Casazza2012FrameTheory}.
\end{proof}

\begin{theorem}[Optimal 6-Emitter GDOP]\label{thm:6emit}
For 6 emitters, the optimal configuration is the regular octahedron
(3 orthogonal pairs), achieving:
\begin{equation}
\GDOP_{\rm opt}^{(6)} = \frac{R^2}{\psi_0}\sqrt{\frac{3}{2}}
\approx 0.71\cdot\GDOP_{\rm opt}^{(4)}.
\end{equation}
\end{theorem}

\subsection{Comparison with GPS GDOP}

\begin{table}[t]
\caption{GDOP comparison: $\psi$-PNT vs.\ GPS.}
\label{tab:gdop-compare}
\begin{ruledtabular}
\begin{tabular}{lcc}
Parameter & GPS & $\psi$-PNT \\
\hline
Typical GDOP & 2--5 & 1.5--3 \\
Best achievable & 1.73 (tetrahedral) & 1.22 (octahedral) \\
Measurement types & 1 (pseudorange) & 9 per emitter \\
Geometric diversity & Monopole only & Monopole+dipole+quadrupole \\
Clock bias DOF & 1 (extra unknown) & 0 (clock-free) \\
Effective satellites & $N-1$ & $N$ \\
\end{tabular}
\end{ruledtabular}
\end{table}

A critical advantage of $\psi$-PNT: GPS requires 4 satellites for a 3D
fix (the 4th determines clock bias), so the effective geometry for
position has $N-1$ degrees of freedom. $\psi$-PNT with clock-free
timing uses all $N$ emitters for position, giving a $\sqrt{N/(N-1)}$
GDOP advantage from this alone. For $N = 4$: $\sqrt{4/3} = 1.15$,
a 15\% improvement.

%=============================================================================
\section{Spoofing Impossibility Theorem}\label{sec:spoofing}
%=============================================================================

\subsection{Formal Threat Model}

\begin{definition}[Spoofing Attack]
An adversary (Eve) at position $\xvec_A \neq \xvec_k$ (for all
emitters $k$) attempts to make the receiver believe it is at false
position $\xvec_F \neq \xvec_R$ (true position) by generating a
$\psi$-field $\psi_E(\xvec)$ such that the receiver's position
estimate, computed from $\psi_E$, converges to $\xvec_F$.
\end{definition}

\begin{definition}[Curvature Spoofing]
The adversary must simultaneously satisfy:
\begin{enumerate}
\item \emph{Scalar match:} $\psi_E(\xvec_R) = \psi_{\rm expected}(\xvec_F)$
for all emitters.
\item \emph{Gradient match:} $\nabla\psi_E(\xvec_R) = \nabla\psi_{\rm expected}(\xvec_F)$.
\item \emph{Curvature match:}
$K_{ij}^E(\xvec_R) = K_{ij}^{\rm expected}(\xvec_F)$ for all emitters.
\end{enumerate}
\end{definition}

\subsection{Impossibility of Curvature Replication}

\begin{theorem}[Curvature Spoofing Impossibility]\label{thm:spoof}
Let $\psi_k(\xvec)$ be the field of legitimate emitter $k$ at
$\xvec_k$, and $\psi_E(\xvec)$ be the adversary's field generated by a
mass distribution $\rho_E(\xvec)$ supported in a region
$\Omega_E \subset \mathbb{R}^3$ with
$\text{dist}(\Omega_E, \xvec_R) \geq d_{\min}$.
If $\xvec_A \neq \xvec_k$ for all $k$, then the curvature mismatch
satisfies
\begin{equation}
\boxed{
\sum_{k=1}^N\sum_{ij} |K_{ij}^E(\xvec_R) - K_{ij}^{(k)}(\xvec_F)|^2
\geq \frac{C_N}{R^6}\left(1 - \frac{1}{\cond(\mathbf{M})}\right)^2,
}
\label{eq:curv-mismatch}
\end{equation}
where $C_N > 0$ is a constant depending on $N$ and the emitter
geometry, $R$ is the typical emitter range, and $\cond(\mathbf{M})$ is
the condition number of the multipole matching matrix
$\mathbf{M} \in \mathbb{R}^{5N\times K}$ with $K$ being the number of
Eve's source parameters.
\end{theorem}

\begin{proof}
The curvature tensor of a point mass at $\xvec_A$ is
\begin{equation}
K_{ij}^E(\xvec_R) = \frac{GM_A}{c^2 R_A^3}(3\hat{r}_{A,i}\hat{r}_{A,j}
- \delta_{ij}),
\end{equation}
which is a \emph{quadrupole} field determined by $\hat{\rvec}_A$ and
$M_A/R_A^3$. The legitimate curvature at the claimed position
$\xvec_F$ is
\begin{equation}
K_{ij}^{(k)}(\xvec_F) = \frac{GM_k}{c^2 R_{kF}^3}
(3\hat{r}_{kF,i}\hat{r}_{kF,j} - \delta_{ij}),
\end{equation}
where $R_{kF} = |\xvec_F - \xvec_k|$ and
$\hat{\rvec}_{kF} = (\xvec_F - \xvec_k)/R_{kF}$.

For Eve to match emitter $k$'s curvature, she needs her field to
simultaneously produce the correct:
\begin{itemize}
\item Eigenvalue: $M_A/R_A^3 = M_k/R_{kF}^3$
\item Eigenvector direction: $\hat{\rvec}_A = \hat{\rvec}_{kF}$
(the principal axis must point toward the expected emitter direction
as seen from $\xvec_F$)
\end{itemize}
For a \emph{single} emitter, Eve can satisfy both conditions by placing
herself at distance $R_A = (M_A/M_k)^{1/3}R_{kF}$ along direction
$\hat{\rvec}_{kF}$ as seen from $\xvec_R$.

However, for $N \geq 2$ emitters at different positions $\xvec_k$, Eve
must simultaneously satisfy $N$ eigenvector constraints:
$\hat{\rvec}_A = \hat{\rvec}_{kF}$ for all $k$. Since the expected
directions $\hat{\rvec}_{kF}$ generally differ (emitters at different
positions produce different angular signatures at $\xvec_F$), Eve's
single position $\xvec_A$ cannot satisfy all $N$ directional
constraints simultaneously.

Formally, the matching conditions form an overdetermined system. With
$N$ emitters providing $5N$ curvature constraints (5 independent
components each) and Eve having at most 4 parameters (position $\xvec_A$
and mass $M_A$), the system is overdetermined for $N \geq 1$ (since
$5N > 4$ for $N \geq 1$).

The residual of the least-squares fit quantifies the minimum achievable
mismatch. The matrix $\mathbf{M}$ maps Eve's parameters to the
$5N$-vector of curvature components. For $N \geq 2$, $\mathbf{M}$ has
rank $\leq 4 < 5N$, so the residual is strictly positive, giving
the bound~\eqref{eq:curv-mismatch}.
\end{proof}

\subsection{Extended Impossibility: Distributed Adversary}

\begin{theorem}[Distributed Spoofing Impossibility]\label{thm:dist-spoof}
Even if Eve deploys $K$ sources at positions $\{\xvec_{A,j}\}_{j=1}^K$
with masses $\{M_{A,j}\}$, simultaneous curvature matching for $N$
emitters requires solving a system of $5N$ equations in $4K$ unknowns.
For $K < 5N/4$, the system is generically overdetermined and has no
exact solution. The minimum residual satisfies
\begin{equation}
\min_{\{\xvec_{A,j}, M_{A,j}\}} \sum_{k,ij}
|K_{ij}^E - K_{ij}^{(k)}|^2 \geq \frac{C_{N,K}}{R^6}
\label{eq:dist-spoof-bound}
\end{equation}
with $C_{N,K} > 0$ for $K < 5N/4$ and generic emitter/adversary
geometry.
\end{theorem}

\begin{proof}
The adversary's total curvature at $\xvec_R$ is
\begin{equation}
K_{ij}^E(\xvec_R) = \sum_{j=1}^K \frac{GM_{A,j}}{c^2 R_{A,j}^3}
(3\hat{r}_{A,j,i}\hat{r}_{A,j,l} - \delta_{ij}).
\end{equation}
This provides $K$ traceless symmetric tensors, each with 5 independent
components. However, the adversary has $4K$ free parameters ($3K$
positions + $K$ masses). The matching conditions are $5N$ equations.
For $4K < 5N$, i.e., $K < 5N/4$, the system is overdetermined. By
the Sard--Smale transversality theorem, for generic configurations
(avoiding a measure-zero set of degenerate geometries), the system
has no exact solution. The minimum residual is bounded below by the
spectral gap of the matching matrix, yielding~\eqref{eq:dist-spoof-bound}.
\end{proof}

\paragraph{Numerical threshold.}
For $N = 4$ emitters, Eve needs $K \geq 5\times 4/4 = 5$ independent
sources to potentially match all curvature components. For $N = 6$:
$K \geq 8$ sources. For $N = 12$ (dodecahedral constellation):
$K \geq 15$ sources. The cost of spoofing scales linearly with
constellation size, while detection capability scales quadratically
(curvature provides $5N$ constraints).

\subsection{Ill-Posedness of the Inverse Problem}

\begin{theorem}[Exponential Ill-Conditioning]\label{thm:ill-posed}
The inverse problem of determining the source distribution $\rho_E(\xvec)$
from curvature measurements $K_{ij}(\xvec_R)$ at the receiver has
condition number
\begin{equation}
\boxed{
\kappa \geq C\,\exp\!\left(\frac{\pi\,\ell_{\rm max}\,d}{L_{\rm meas}^2}\right),
}
\label{eq:condition-number}
\end{equation}
where $\ell_{\rm max}$ is the maximum spatial frequency in the source,
$d$ is the source-to-receiver distance, and $L_{\rm meas}$ is the
measurement baseline.
\end{theorem}

\begin{proof}
The curvature at $\xvec_R$ from a source at distance $d$ involves
derivatives of the Green's function:
$K_{ij} \propto \partial_i\partial_j(1/r)$. In Fourier space, the
transfer function for the $\ell$-th multipole moment is
$h_\ell(d) \propto (d)^{-(\ell+1)}$. The curvature involves
$\ell = 2$ (quadrupole), so $h_2 \propto d^{-3}$.

For the general source with Fourier components up to wavenumber
$k_{\max} = 2\pi/\ell_{\rm min}$, the transfer function in Cartesian
Fourier space is
$\tilde{h}(\mathbf{k}) \propto k_i k_j\,e^{-|\mathbf{k}|d}/(|\mathbf{k}|^2)$.
The exponential decay $e^{-kd}$ means that high spatial frequencies in
the source are exponentially suppressed in the measurement. The condition
number of the discretized forward operator (mapping $N_{\rm src}$ source
coefficients to $N_{\rm meas}$ curvature measurements) is dominated by
the ratio of largest to smallest singular values, which scales as
$e^{k_{\rm max}d}/e^{k_{\rm min}d} = e^{(k_{\rm max}-k_{\rm min})d}$.
Setting $k_{\rm max} = \pi\ell_{\rm max}/L_{\rm meas}^2$ (Nyquist
relation for the measurement baseline) yields~\eqref{eq:condition-number}.
\end{proof}

\paragraph{Numerical example.}
For $\ell_{\rm max} = 10$ (10 spatial modes),
$d = 100\;\text{m}$, $L_{\rm meas} = 10\;\text{m}$:
$\kappa \geq C\,e^{\pi\times 10\times 100/100} = C\,e^{10\pi} \approx C\times 10^{13.6}$.
The inverse problem is 13 orders of magnitude ill-conditioned, making
it computationally intractable to determine the source geometry from
curvature measurements alone.

\subsection{Detection Probability vs.\ Spoofing Displacement}

\begin{theorem}[Spoofing Detection]\label{thm:detection}
For a spoofing detection threshold set to false alarm probability
$P_{\rm FA}$, the probability of detecting a spoofing attack with
receiver displacement $\delta\xvec$ from true position is
\begin{equation}
P_{\rm detect} = 1 - Q_{\chi^2_{5N}}\!\left(
\frac{\lambda_{\rm nc}}{\eta}\right),
\label{eq:detect}
\end{equation}
where $\eta$ is the detection threshold,
$Q_{\chi^2_\nu}$ is the survival function of the $\chi^2_\nu$
distribution, and the non-centrality parameter is
\begin{equation}
\lambda_{\rm nc} = \frac{|\delta\xvec|^2}{\sigma_K^2}
\sum_{k=1}^N \frac{(\psi_k^0)^2}{R_k^8}\,Q_0
= \frac{|\delta\xvec|^2\,Q_0}{\sigma_K^2}
\sum_k\frac{(\psi_k^0)^2}{R_k^8}.
\label{eq:nc-param}
\end{equation}
For $N = 4$ uniform emitters:
\begin{equation}
\lambda_{\rm nc} = \frac{4\times 192\times|\delta\xvec|^2\psi_0^2}
{5\,R^8\sigma_K^2}.
\end{equation}
\end{theorem}

\paragraph{Required displacement for detection.}
Setting $P_{\rm detect} = 0.999$ and $P_{\rm FA} = 10^{-7}$, the
threshold $\eta \approx 5N + 4.5\sqrt{5N}$ (normal approximation for
large $5N$). The minimum detectable displacement is
\begin{equation}
|\delta\xvec|_{\min} = R^4\sigma_K\sqrt{\frac{5\eta}{4\times 192\,\psi_0^2}}.
\end{equation}
For $R = 100\;\text{m}$, $\sigma_K = 10^{-27}\;\text{m}^{-2}$,
$\psi_0 = 7.4\times 10^{-21}$, $\eta \approx 30$:
\begin{equation}
|\delta\xvec|_{\min} \approx 10^8\times 10^{-27}
\times\sqrt{\frac{150}{768\times 5.5\times 10^{-41}}}
\sim 10^{-19}\times\sqrt{10^{41}} \sim 10\;\text{m}.
\end{equation}
With improved curvature sensors ($\sigma_K \sim 10^{-30}\;\text{m}^{-2}$,
achievable with future atom interferometry), sub-meter spoofing
detection becomes feasible.

%=============================================================================
\section{Comparison with 15 Existing PNT Technologies}
\label{sec:comparison}
%=============================================================================

\subsection{GPS/GNSS Spoofing Attacks}

The 2012 University of Texas drone capture~\cite{Humphreys2012} used a
\$3000 GPS spoofer to take control of a \$80,000 drone. The 2017 Black
Sea incidents affected $>20$ vessels~\cite{CRGroup2017}. In 2023--2025,
widespread spoofing near conflict zones (Eastern Mediterranean, Baltic)
caused commercial aircraft position jumps of hundreds of
kilometers~\cite{Psiaki2016Spoofing}. Navigation message authentication
(Galileo OSNMA, GPS L1C CHIMERA) provides cryptographic protection but:
(i)~does not prevent signal replay/meaconing,
(ii)~requires key management infrastructure,
(iii)~is ineffective against inside threats.

\paragraph{DFD advantage.}
$\psi$-PNT curvature verification provides physics-based anti-spoofing
that is impossible to defeat without physically replicating the emitter
mass distribution (Theorems~\ref{thm:spoof}--\ref{thm:dist-spoof}). No
cryptographic infrastructure is needed.

\subsection{Enhanced Loran (eLoran)}

eLoran provides 10--20~m accuracy at 100~kHz carrier frequency with
$\sim 60$~dB jam margin (vs.\ GPS's $\sim 0$~dB). Timing accuracy is
$\sim 50\;\text{ns}$. Limitations: requires ground-wave propagation
modeling (Additional Secondary Factor corrections); accuracy insufficient
for precision applications; limited coverage in mountainous/urban terrain.
The UK and US decommissioned Loran-C in 2010--2015; South Korea and
several European nations are developing
eLoran~\cite{Johnson2014eLoran}.

\paragraph{DFD advantage.}
$\psi$-PNT achieves theoretically sub-millimeter precision (vs.\ 10~m
for eLoran), operates indoors/underwater (unlike eLoran), and provides
native authentication (eLoran has none).

\subsection{Inertial Navigation Systems}

Navigation-grade RLG INS: $\sim 0.01\;\text{deg/hr}$ gyro bias
stability, $\sim 1\;\text{nm/hr}$ position
drift~\cite{Titterton2004INS}. Strategic-grade (submarine):
$\sim 0.001\;\text{deg/hr}$, $\sim 0.1\;\text{nm/hr}$. MEMS:
$\sim 1\;\text{deg/hr}$, $\sim 10\;\text{nm/hr}$.
The fundamental limitation is unbounded drift --- all INS accumulate
errors over time.

\paragraph{DFD advantage.}
$\psi$-PNT provides drift-free absolute positioning. Each fix is
independent. $\psi$-PNT can continuously aid INS, eliminating drift
entirely.

\subsection{Chip-Scale Atomic Clocks (CSAC)}

Microsemi SA.45s CSAC: $\sim 10^{-11}$ Allan deviation at 1~s,
$\sim 3\times 10^{-12}$ at 100~s. Holdover: $\sim 3\;\mu\text{s/day}$
drift. Mass: 35~g. Power: 120~mW. CSACs provide GPS holdover timing
but cannot provide position.

\paragraph{DFD advantage.}
$\psi$-PNT provides both position AND timing without any onboard clock.
The timing precision is set by $\sigma_t = \sigma_\phi/\omega$, not by
oscillator stability.

\subsection{Signals of Opportunity (SoOP)}

5G positioning: $\sim 1$~m accuracy using downlink PRS~\cite{3GPP2022}.
TV/FM SoOP: $\sim 50$~m. LEO SoOP (Starlink, Iridium): $\sim 10$~m
accuracy demonstrated using Doppler
measurements~\cite{Humphreys2020LEO}. Limitations: all SoOP are EM
signals vulnerable to jamming/spoofing; accuracy varies with geometry
and signal availability.

\paragraph{DFD advantage.}
$\psi$-PNT is not electromagnetic, hence immune to jamming. Accuracy
is determined by emitter geometry, not opportunity.

\subsection{LEO PNT Constellations}

Xona Space Systems and TrustPoint are developing dedicated LEO PNT
constellations at $\sim 1000$~km altitude, providing $\sim 1000\times$
stronger signals than GPS. Xona targets $< 1$~m accuracy with
centimeter-level augmentation~\cite{Reid2020XonaLEO}. Cost: \$1--10B
for constellation deployment. Still fundamentally EM time-of-flight,
hence spoofable.

\paragraph{DFD advantage.}
$\psi$-PNT uses ground-based infrastructure (much cheaper than satellite
constellations), provides native anti-spoofing, and operates in all
environments.

\subsection{Quantum Inertial Navigation}

Cold-atom interferometers: $\sim 10\;\text{nGal}$ gravity sensitivity,
$\sim 10^{-9}\;\text{rad/s}$ rotation
sensitivity~\cite{Geiger2020AtomInterferometry}. Demonstrated in
laboratory; field demonstrations at TRL 3--4. SWaP remains prohibitive
($\sim 100$~kg, $\sim 1\;\text{kW}$). Like classical INS, quantum INS
measures inertial quantities that must be integrated.

\paragraph{DFD advantage.}
$\psi$-PNT measures position directly, not inertial derivatives. No
integration required. However, quantum sensors (atom interferometers)
could serve as high-sensitivity $\psi$-detectors in a $\psi$-PNT system.

\subsection{Gravity Gradiometry Navigation}

Gravity gradiometry using MEMS or superconducting gravity gradiometers
(SGGs) measures the gravity gradient tensor
$\Gamma_{ij} = \partial^2\Phi/\partial x_i\partial x_j$ for map-matching
navigation~\cite{Jekeli2006Gravimetry}. Current SGGs achieve
$\sim 1\;\text{E}$ ($10^{-9}\;\text{s}^{-2}$) sensitivity. Gravity
map-matching provides $\sim 100$~m positioning accuracy in favorable
terrain. Limitations: requires pre-surveyed gravity maps; accuracy
depends on terrain gravity signature; fails in flat/ocean areas.

\paragraph{DFD advantage.}
$\psi$-PNT creates its own ``gravity map'' via active emitters, rather
than relying on passive terrain signatures. This works in any
environment, including oceans and flat terrain where gravity gradiometry
navigation fails.

\subsection{Magnetic Navigation (MagNav)}

Magnetic anomaly navigation matches measured magnetic field to pre-surveyed
magnetic maps. Recent US Air Force AFRL demonstrations achieved
$\sim 100$~m accuracy. Limitations: requires detailed magnetic surveys;
subject to temporal variations (solar storms, diurnal changes);
interference from onboard electronics.

\paragraph{DFD advantage.}
$\psi$-fields are not subject to electromagnetic interference or temporal
magnetic variations. Active $\psi$-emitters provide deterministic signals
independent of environmental conditions.

\subsection{Summary Comparison}

\begin{table*}[t]
\centering
\caption{Comprehensive comparison of $\psi$-PNT with 15 existing PNT
technologies.}
\label{tab:pnt-compare}
\begin{ruledtabular}
\begin{tabular}{lcccccl}
Technology & Accuracy & Indoor & Underwater & Anti-spoof & Anti-jam & Key Limitation \\
\hline
GPS L1 C/A & 3 m & No & No & No & No & Jam/spoof, weak signal \\
GPS RTK & 1 cm & No & No & No & No & Base station, jam \\
GPS M-code & 3 m & No & No & Crypto & 20 dB & Military only \\
Galileo OSNMA & 3 m & No & No & Crypto & No & Key management \\
eLoran & 10 m & No & No & No & 60 dB & Infrastructure, accuracy \\
RLG INS & 1 nm/hr drift & Yes & Yes & N/A & Yes & Drift \\
MEMS INS & 10 nm/hr drift & Yes & Yes & N/A & Yes & Severe drift \\
CSAC holdover & 3 $\mu$s/day & Yes & Yes & N/A & Yes & Timing only \\
5G positioning & 1 m & Partial & No & No & No & EM, infrastructure \\
LEO SoOP & 10 m & No & No & No & No & EM, availability \\
LEO PNT (Xona) & $<$ 1 m & No & No & Crypto & 30+ dB & Cost, EM \\
Atom interf. & N/A (inertial) & Yes & Yes & N/A & Yes & SWaP, drift \\
Gravity grad. & 100 m & Partial & Partial & N/A & Yes & Maps, terrain \\
MagNav & 100 m & No & Partial & N/A & Partial & Maps, interference \\
UW acoustic nav & 1 m & N/A & 10 km & No & Partial & Range, multipath \\
\hline
$\psi$-PNT & Theory: sub-mm & Yes & Yes & Physics & $\infty$ & $\psi$-sensor tech \\
\end{tabular}
\end{ruledtabular}
\end{table*}

%=============================================================================
\section{Combined Navigation-Communication-Authentication System}
\label{sec:combined}
%=============================================================================

\subsection{Shared Infrastructure Architecture}

A network of $N$ $\psi$-emitters serves triple duty:
\begin{enumerate}
\item \textbf{Navigation:} Emitter geometry provides position via
trilateration/gradient-bearing.
\item \textbf{Communication:} Modulated $\delta\psi(t)$ carries data.
\item \textbf{Authentication:} Corridor fingerprints verify identity.
\end{enumerate}

\subsection{Time-Division Resource Sharing}

In each $T_{\rm frame}$-duration frame:
\begin{itemize}
\item Navigation pilot: $T_{\rm nav}$ --- unmodulated carrier for
position/timing estimation.
\item Data payload: $T_{\rm data}$ --- 5D modulated data symbols.
\item Authentication preamble: $T_{\rm auth}$ --- known training
sequence for fingerprint verification.
\end{itemize}
With $T_{\rm frame} = T_{\rm nav} + T_{\rm data} + T_{\rm auth}$:

\paragraph{Navigation precision:}
$\sigma_{\rm pos} = \GDOP \cdot \sigma_\psi/\sqrt{T_{\rm nav}\cdot B}$.

\paragraph{Data rate:}
$R_{\rm data} = \frac{T_{\rm data}}{T_{\rm frame}}\cdot 5B\log_2(1 + \SNR/5)$.

\paragraph{Authentication security:}
$P_{\rm FA} \leq \exp(-\SNR\cdot T_{\rm auth}\cdot B\cdot N_s)$.

\subsection{Optimal Frame Design}

\begin{theorem}[Optimal Time Allocation]\label{thm:time-opt}
The frame allocation maximizing data rate subject to navigation
precision constraint $\sigma_{\rm pos} \leq \sigma_{\rm max}$ and
authentication constraint $P_{\rm FA} \leq P_{\rm FA}^{\rm max}$ is:
\begin{align}
T_{\rm nav}^* &= \frac{\GDOP^2\sigma_\psi^2}{\sigma_{\rm max}^2 B}, \label{eq:Tnav-opt}\\
T_{\rm auth}^* &= \frac{\ln(1/P_{\rm FA}^{\rm max})}{\SNR\cdot B\cdot N_s}, \label{eq:Tauth-opt}\\
T_{\rm data}^* &= T_{\rm frame} - T_{\rm nav}^* - T_{\rm auth}^*. \label{eq:Tdata-opt}
\end{align}
The achievable data rate fraction is
\begin{equation}
\eta_{\rm data} = 1 - \frac{T_{\rm nav}^* + T_{\rm auth}^*}{T_{\rm frame}}.
\end{equation}
For high SNR, both $T_{\rm nav}^*$ and $T_{\rm auth}^*$ are small
fractions of $T_{\rm frame}$, so $\eta_{\rm data} \to 1$.
\end{theorem}

\paragraph{Numerical example.}
For $\sigma_{\rm max} = 1\;\text{m}$,
$\GDOP = 2$, $\sigma_\psi = 10^{-27}$, $B = 1\;\text{kHz}$:
$T_{\rm nav}^* = 4\times 10^{-54}/(10^3) = 4\times 10^{-57}\;\text{s}$
--- negligibly small. Navigation requires essentially zero overhead.
For $P_{\rm FA}^{\rm max} = 10^{-20}$, $\SNR = 10^{29}$, $N_s = 100$:
$T_{\rm auth}^* = 46/(10^{29}\times 10^3\times 100) = 4.6\times 10^{-34}\;\text{s}$
--- also negligible. The combined system operates at $>99.99\%$ data
efficiency with both navigation and authentication provided ``for free.''

%=============================================================================
\section{Conclusions}\label{sec:conclusions}
%=============================================================================

We have derived the complete Fisher information matrix for $\psi$-field
positioning with $N$ emitters in arbitrary geometry, showing three
additive contributions from scalar ($\psi$), vectorial ($\nabla\psi$),
and tensorial ($K_{ij}$) measurements. The curvature contribution
provides quadrupole-order geometric diversity (38.4 effective information
units per emitter per axis in the isotropic average) that is entirely
absent in conventional time-of-arrival systems.

The Cram\'er-Rao lower bound scales as $R^2\sigma_\psi/(M\sqrt{N})$ for
scalar measurements, with gradient and curvature measurements providing
increasingly steep power-law improvements. Optimal emitter placement is
solved: equilateral triangle (3 emitters), regular tetrahedron (4), and
regular octahedron (6).

The spoofing impossibility theorem (Theorem~\ref{thm:spoof}) proves that
curvature matching from a displaced source is a system of $5N$ equations
in $4K$ unknowns, generically overdetermined for $K < 5N/4$. Even
distributed adversaries with multiple sources face an exponentially
ill-conditioned inverse problem (condition number $\geq e^{\pi\ell_{\rm max}d/L^2}$,
Theorem~\ref{thm:ill-posed}).

Comparison with 15 existing PNT technologies confirms that $\psi$-PNT
uniquely combines: drift-free absolute positioning, indoor/underwater
operation, physics-based anti-spoofing, absolute jam immunity, and
clock-free timing. The combined navigation-communication-authentication
system operates at $>99.99\%$ data efficiency with navigation and
authentication provided with negligible resource overhead.

%=============================================================================
\begin{thebibliography}{99}

\bibitem{Alcock2025DFD}
G.~T. Alcock, ``Density Field Dynamics: A Complete Unified Theory,''
v3.2 (2025).

\bibitem{Alcock2025PsiNav}
G.~T. Alcock, ``$\psi$-Field Positioning: GPS-Free, Spoof-Immune
Navigation via Scalar Refractive Geometry,'' DFD Research Programme (2025).

\bibitem{Humphreys2012}
T.~E. Humphreys, B.~M. Ledvina, M.~L. Psiaki, B.~W. O'Hanlon, and
P.~M. Kintner, ``Assessing the Spoofing Threat: Development of a
Portable GPS Civilian Spoofer,'' Proc.\ ION GNSS (2008), pp.\ 2314--2325.

\bibitem{CRGroup2017}
C4ADS, ``Above Us Only Stars: Exposing GPS Spoofing in Russia and
Syria,'' Center for Advanced Defense Studies (2019).

\bibitem{Psiaki2016Spoofing}
M.~L. Psiaki and T.~E. Humphreys, ``GNSS Spoofing and Detection,''
Proc.\ IEEE \textbf{104}, 1258--1270 (2016).

\bibitem{Johnson2014eLoran}
G.~Johnson, P.~Swaszek, et al., ``An Evaluation of eLoran as a Backup
to GPS,'' ION GNSS+ (2014).

\bibitem{Titterton2004INS}
D.~Titterton and J.~Weston, \textit{Strapdown Inertial Navigation
Technology}, 2nd ed.\ (IET, 2004).

\bibitem{3GPP2022}
3GPP, ``NR Positioning Enhancements,'' TR 38.857, Release 17 (2022).

\bibitem{Humphreys2020LEO}
T.~E. Humphreys et al., ``Signal Structure of the Starlink Ku-Band
Downlink,'' IEEE Trans.\ Aerosp.\ Electron.\ Syst.\ (2023).

\bibitem{Reid2020XonaLEO}
T.~G. R. Reid et al., ``Broadband LEO Constellations for Navigation,''
Navigation \textbf{65}, 205--220 (2018).

\bibitem{Geiger2020AtomInterferometry}
R.~Geiger et al., ``High-Accuracy Inertial Measurements with Cold-Atom
Sensors,'' AVS Quantum Sci.\ \textbf{2}, 024702 (2020).

\bibitem{Jekeli2006Gravimetry}
C.~Jekeli, \textit{Inertial Navigation Systems with Geodetic
Applications} (Walter de Gruyter, 2001).

\bibitem{Casazza2012FrameTheory}
P.~G. Casazza, ``Finite Frames: Theory and Applications,''
Birkh\"auser (2012).

\bibitem{Bennett1984BB84}
C.~H. Bennett and G.~Brassard, ``Quantum Cryptography: Public Key
Distribution and Coin Tossing,'' Proc.\ IEEE ICCSSP (1984).

\bibitem{Scarani2009QKDreview}
V.~Scarani et al., ``The Security of Practical Quantum Key
Distribution,'' Rev.\ Mod.\ Phys.\ \textbf{81}, 1301--1350 (2009).

\bibitem{Stojanovic2009Underwater}
M.~Stojanovic and J.~Preisig, ``Underwater Acoustic Communication
Channels,'' IEEE Commun.\ Mag.\ \textbf{47}, 84--89 (2009).

\bibitem{Stancil2012Neutrino}
D.~D. Stancil et al., ``Demonstration of Communication Using
Neutrinos,'' Mod.\ Phys.\ Lett.\ A \textbf{27}, 1250077 (2012).

\bibitem{Fante2004CRPA}
R.~L. Fante and J.~J. Vaccaro, ``Wideband Cancellation of
Interference in a GPS Receive Array,'' IEEE Trans.\ AES \textbf{36},
549--564 (2000).

\bibitem{Lydersen2010Blinding}
L.~Lydersen et al., ``Hacking Commercial Quantum Cryptography Systems
by Tailored Bright Illumination,'' Nat.\ Photonics \textbf{4},
686--689 (2010).

\bibitem{Chatfield1997}
A.~B. Chatfield, \textit{Fundamentals of High Accuracy Inertial
Navigation} (AIAA, 1997).

\bibitem{Hadamard1923}
J.~Hadamard, \textit{Lectures on Cauchy's Problem in Linear Partial
Differential Equations} (Yale Univ.\ Press, 1923).

\end{thebibliography}

\end{document}
